% Suggested insertion point:
% Chapter 5, after the simulation results or before the final discussion of failure cases.

\subsection{Real-Robot Validation}
\label{subsec:real_robot_validation}

The main quantitative benchmark in this thesis is conducted in simulation in order to evaluate
many task instances under controlled and reproducible conditions. To complement this
benchmark, an additional real-robot validation was performed on the physical UR5 tabletop
setup. The purpose of this experiment is not to replace the large-scale simulated tier
evaluation, but to demonstrate that the proposed perception-planning-execution pipeline can
be transferred to a physical robot and to identify practical failure modes caused by perception
noise, calibration error, grasping uncertainty, and contact-sensitive LEGO insertion.

The real-robot tasks are organized according to the same increasing-complexity principle as
the simulated tiers, but the physical validation is restricted to three tiers. Table~\ref{tab:real_robot_tiers}
summarizes the hardware validation tasks. Tier~1 isolates the single-step perception-to-execution
loop. Tier~2 evaluates vertical multi-step stacking and therefore tests error accumulation over
several placements. Tier~3 uses a bridge structure with three identical white 2x4 bricks and then
adds a red 2x4 brick and a yellow 2x2 brick. This tier stresses repeated-object handling,
color-sensitive perception, mixed part sizes, longer action horizons, and taller placement.

\begin{table}[t]
    \centering
    \caption{Real-robot validation tasks grouped by increasing structural complexity.}
    \label{tab:real_robot_tiers}
    \begin{tabular}{llll}
        \toprule
        Tier & Product & Main difficulty & Planned steps \\
        \midrule
        1 & Single yellow 2x4 & Single-object pose estimation and placement & 1 \\
        2 & Three 2x4 vertical stack & Multi-step stacking and height accumulation & 3 \\
        3 & Bridge + red 2x4 + yellow 2x2 & Repeated identical parts, mixed colors, taller structure & 5 \\
        \bottomrule
    \end{tabular}
\end{table}

\paragraph{Protocol and recorded metrics.}
For each tier, the system starts from a fixed observation pose and receives the corresponding
product YAML file. The planner first generates the assembly order. The vision module then
estimates the target 6D pose and writes the result to a structured YAML file. The execution
module converts this result into a robot-specific pick-and-place task and executes the
corresponding URScript motion sequence.

Each trial records task-level success and end-to-end runtime. A trial is counted as successful
only when all planned actions are completed and the final structure is assembled correctly.
The reported runtime is the total wall-clock time of the complete pipeline and includes
planning, visual perception, 6D pose estimation, YAML conversion, communication overhead,
and physical robot execution.

\begin{table}[t]
    \centering
    \caption{Real-robot validation results. Runtime is measured as end-to-end wall-clock time
    for the full planning--perception--execution pipeline.}
    \label{tab:real_robot_results}
    \begin{tabular}{lrrrrrr}
        \toprule
        Tier & Trials & Steps/trial & Successful trials & Success rate & Mean total time & Assembly success \\
        \midrule
        1 & 10 & 1 & 9/10 & 90.0\% & 55.78\,s & 9/10 \\
        2 & 10 & 3 & 9/10 & 90.0\% & 167.98\,s & 9/10 \\
        3 & 10 & 5 & 8/10 & 80.0\% & 183.33\,s & 8/10 \\
        \bottomrule
    \end{tabular}
\end{table}

\begin{table}[t]
    \centering
    \caption{End-to-end real-robot runtime statistics over ten trials per tier. Times are
    reported in seconds and include planning, perception, conversion, communication, and
    robot execution.}
    \label{tab:real_robot_timing}
    \begin{tabular}{lrrrrr}
        \toprule
        Tier & Mean & Std. dev. & Median & Min. & Max. \\
        \midrule
        1 & 55.78 & 4.29 & 57.90 & 45.34 & 59.35 \\
        2 & 167.98 & 8.96 & 170.23 & 152.34 & 177.34 \\
        3 & 183.33 & 7.38 & 182.30 & 170.35 & 195.32 \\
        \bottomrule
    \end{tabular}
\end{table}

\paragraph{Analysis.}
The real-robot results should be interpreted as a hardware validation of the integrated system
rather than as a large-scale statistical benchmark. Tier~1 and Tier~2 both complete 9 out of
10 trials successfully, while Tier~3 completes 8 out of 10 trials. The decrease in Tier~3 is
expected because it combines repeated identical white bricks, color-sensitive perception, a
2x2 part, five physical actions, and a taller final product. The end-to-end runtime increases
from 55.78 seconds in Tier~1 to 167.98 seconds in Tier~2 and 183.33 seconds in Tier~3,
reflecting the longer action horizon and additional perception/execution steps.

The asynchronous perception-execution schedule should be discussed separately from the
success-rate results. Its purpose is not to improve pose accuracy directly, but to reduce
pipeline idle time. Therefore, the relevant evidence is whether the next vision output is often
ready before the current robot action has completed. If so, the sequential perception wait can
be partially hidden behind physical execution time. This supports the broader claim that the
pipeline can be engineered as an integrated robotic system rather than as isolated perception,
planning, and execution modules.
